\documentclass[../../../include/open-logic-section]{subfiles}

\begin{document}

\olfileid{sth}{card-arithmetic}{opps}
\olsection{Netepaké Operasi Dhasar}

Amarga ora perlu nggatekaké urutan, aritmetika kardinal luwih
gampang ditetepaké tinimbang aritmetika ordinal. Kita bakal
netepaké panambahan, ping-pingan, lan pemangkatan bebarengan.

\begin{defn}
Yèn $\cardfont{a}$ lan $\cardfont{b}$ iku kardinal:
\begin{align*}
	\cardfont{a} \cardplus \cardfont{b} &\defis
	\card{\cardfont{a} \disjointsum \cardfont{b}}\\
	\cardfont{a} \cardtimes \cardfont{b} &\defis
	\card{\cardfont{a} \times \cardfont{b}}\\
	\cardexpo{\cardfont{a}}{\cardfont{b}} &\defis
	\card{\funfromto{\cardfont{b}}{\cardfont{a}}}
\end{align*}
ing kéné $\funfromto{X}{Y} = \Setabs{f}{f \text{ iku fungsi } X \to Y}$.
(Gampang dituduhaké yèn $\funfromto{X}{Y}$ ana kanggo
sembarang himpunan $X$ lan $Y$; iki dadi latihan.)
\end{defn}

\begin{prob}
Buktèkna ing $\Zminus$ yèn $\funfromto{X}{Y}$ ana kanggo
sembarang himpunan $X$ lan~$Y$. Kanthi makarya ing $\ZF$,
itung $\setrank{\funfromto{X}{Y}}$ saka
$\setrank{X}$ lan $\setrank{Y}$, kaya ing
\olref[sth][ord-arithmetic][using-addition]{rankcomputation}.
\end{prob}

Definisi iki perlu diterangaké sedhela. Kanggo panambahan:
kita nganggo gagasan gunggung pisah, $\disjointsum$,
kaya kang ditetepaké ing
\olref[ord-arithmetic][add]{defdissum}; gampang dideleng
yèn definisi iki menehi asil kang trep kanggo kasus winates.
Kanggo ping-pingan:
\olref[sfr][set][pai]{cardnmprod} kandha yèn $A$
duwé $n$ unsur lan $B$ duwé $m$ unsur, mula $A \times B$
duwé $n \cdot m$ unsur; dadi definisi kita nggedhèkaké gagasan
iki menyang ping-pingan transfinit. Pemangkatan uga padha:
kita nggedhèkaké gagasan saka kasus winates menyang
kasus transfinit. Pancen, ing sawetara babagan,
aritmetika kardinal transfinit luwih mirip aritmetika
“biasa” tinimbang aritmetika ordinal transfinit:

\begin{prop}\ollabel{cardplustimescommute}
$\cardplus$ lan $\cardtimes$ iku komutatif lan asosiatif.
\end{prop}

\begin{proof}
Kanggo komutativitas, miturut
\olref[cardinals][cardsasords]{lem:CardinalsBehaveRight}
cukup wigatèkna yèn $\cardeq{(\cardfont{a} \disjointsum
\cardfont{b})}{(\cardfont{b} \disjointsum \cardfont{a})}$ lan
$\cardeq{(\cardfont{a} \times \cardfont{b})}{(\cardfont{b} \times
\cardfont{a})}$. Asosiativitas kita pasrahaké minangka latihan.
\end{proof}

\begin{prob}
Buktèkna yèn $\cardplus$ lan $\cardtimes$ asosiatif.
\end{prob}

\begin{prop}
$A$ tanpa wates yen lan mung yen $\card{A} \cardplus 1 = 1 \cardplus \card{A} = \card{A}$.
\end{prop}

\begin{proof}
Kaya ing
\olref[cardinals][classing]{generalinfinitycharacter}, saka
\olref[ord-arithmetic][using-addition]{ordinfinitycharacter} lan
\olref[cardinals][cardsasords]{lem:CardinalsBehaveRight}.
\end{proof}

Iki nerangaké sebabé kita kudu nganggo simbol béda kanggo
panambahan lan ping-pingan ordinal lan kardinal: operasi-operasi
iki pancèn \emph{béda}. Rong asil sabanjuré nuduhaké yèn
pemangkatan ordinal lan kardinal uga béda.
(Élinga yèn \olref[z][infinity-again]{defnomega}
ngakibataké $2 = \{0,
1\}$):

\begin{lem}\ollabel{lem:SizePowerset2Exp}
$\card{\Pow{A}} = \cardexpo{2}{\card{A}}$, kanggo sembarang $A$.
\end{lem}

\begin{proof}
Kanggo saben subhimpunan $B \subseteq A$, tetepna $\chi_B \in \funfromto{A}{2}$ kanthi:
\begin{align*}
	\chi_{B}(x) &\defis
	\begin{cases}
		1 & \text{yèn }x\in B\\
		0 & \text{yèn ora.}
	\end{cases}
\end{align*}
Saiki tetepna $f(B) = \chi_B$; iki netepaké !!a{bijection} $f \colon \Pow{A}
\to \funfromto{A}{2}$. Dadi $\cardeq{\Pow{A}}{\funfromto{A}{2}}$. Mula
$\cardeq{\Pow{A}}{\funfromto{\card{A}}{2}}$, saéngga
$\card{\Pow{A}} = \card{\funfromto{\card{A}}{2}} =
2^{\card{A}}$.
\end{proof}

Bukti ringkes iki sakjané nyakup rembugan ing
\olref[sfr][siz][red-alt]{sec}. Ing kono kita nuduhaké carané
“nyuda” sifat ora bisa diétung saka $\Pow{\omega}$
menyang sifat ora bisa diétung saka himpunan untai biner
tanpa wates, $\Bin^\omega$. Ing intiné, $\Bin^{\omega}$
iku $\funfromto{\omega}{2}$; lan bukti sadurungé nuduhaké
yèn panalaran ing \olref[sfr][siz][red-alt]{sec} uga lumaku
yèn sembarang himpunan~$A$ nggantèkaké~$\omega$.
Asil iku uga menehi pratelan cekak babagan pemangkatan
kardinal:

\begin{cor}\ollabel{cantorcor}
$\cardfont{a} < \cardexpo{2}{\cardfont{a}}$ kanggo sembarang kardinal~$\cardfont{a}$.
\end{cor}

\begin{proof}
Saka Teorema Cantor (\olref[sfr][siz][car]{thm:cantor}) lan
\olref{lem:SizePowerset2Exp}.
\end{proof}
\noindent
Dadi $\omega < \cardexpo{2}{\omega}$. Nanging wigatèkna:
iki asil babagan pemangkatan \emph{kardinal}. Bandhingna
karo pemangkatan \emph{ordinal}, amarga ing kasus iku
$\omega = \ordexpo{2}{\omega}$
(delengen \olref[ord-arithmetic][expo]{sec}).

Isih babagan pemangkatan kardinal, kita bisa nerangaké
kanthi luwih trep “cara” $\Real$ iku
!!{nonenumerable}.

\begin{thm}\ollabel{continuumis2aleph0}
$\card{\Real} = \cardexpo{2}{\omega}$
\end{thm}

\begin{proof}[Rangka bukti]
Ana akèh cara mbuktèkaké iki. Cara kang paling langsung
yaiku nuduhaké $\cardle{\Pow{\omega}}{\Real}$ lan
$\cardle{\Real}{\Pow{\omega}}$, banjur nganggo
Schr\"oder-Bernstein kanggo nyimpulaké
$\cardeq{\Real}{\Pow{\omega}}$, lan
\olref[card-arithmetic][opps]{lem:SizePowerset2Exp} kanggo
nyimpulaké $\card{\Real} = \cardexpo{2}{\omega}$.
Kita pasrahaké minangka latihan kang maringi pangerten
kanggo netepaké injeksi $f \colon
\Pow{\omega} \to \Real$ lan $g \colon \Real \to \Pow{\omega}$.
\end{proof}

\begin{prob}
Rampungna bukti
\olref[sth][card-arithmetic][opps]{continuumis2aleph0}, kanthi
nuduhaké $\cardle{\Pow{\omega}}{\Real}$ lan
$\cardle{\Real}{\Pow{\omega}}$.
\end{prob}

\end{document}
